\section{Exercise 1}
The Poission distribution is given by 
\begin{equation}\label{eq:poisson}
    P_{\lambda}(k) = \frac{\lambda^k e^{-\lambda}}{k!}
\end{equation}
with $k$ integer and the mean $\lambda$ is a positive number. The factorial of k will produce very large floating point numbers which causes an \texttt{Overflow} error. Therefore, we can also calculate the log of $P_{\lambda}(k)$, which can be written as (according to WolframAlpha)
\begin{equation}
     k\log(\lambda) - \lambda - \log(\Gamma(1 + k))
\end{equation}
where $\Gamma(1+k) = (1+k-1)! = k!$. We can approximate the Gamma function by the Stirling approximation, 
\begin{equation}
    \Gamma(1+k) = k! \approx \sqrt{2\pi k}\left(\frac{k}{e}\right)^k.
\end{equation}

However, this still results in an overflow for larger values of $k$. Therefore, we switched approach and did not use the Gamma function. Instead, we used Ramanujan's approximation for $\log(k!)$:
\begin{equation}\label{eq:ram}
    \log(k!) \approx k\log(k) - k + \frac{\log(k(1+4k(1 + 2k)))}{6} + \frac{log(\pi)}{2}
\end{equation}

An alternate form of the log of $P_{\lambda}(k)$ is 
\begin{equation}
    k\log(\lambda) - \lambda - \log(k!)
\end{equation}

Combining this with \autoref{eq:ram} gives:
\begin{equation}\label{eq:logP}
    \log(P_{\lambda}(k)) = k\log(\lambda) - \lambda - k\log(k) + k - \frac{\log(k(1+4k(1 + 2k)))}{6} - \frac{log(\pi)}{2}
\end{equation}

In the code, we first try to use \autoref{eq:poisson} as this does not overflow for small $k$ values. If it overflows, we use \autoref{eq:logP} and simply return $\exp(\log(P_{\lambda}(k)))$.
NumPy's \texttt{exp()} function can also overflow, and instead of raising an error it will return zero or infinity. In this case, the value of $P$ just contains too many digits, and the only way to get a reasonable output is to use the Decimal module, which can handle these sitautions. However, this only happens in extreme cases, and we put this as an extra in code.  

We make sure that all $k$ and $\lambda$ values we pass to the Poisson distribution and factorial function are \texttt{float32}s. In the factorial function, we round $k$ to make it an integer.

See the code below for our implementation of this exercise:
\lstinputlisting{NUR_Handin_1_1.py}

As an additional test of the robustness of our algorithm, we also tested 10000 combinations of values between $\lambda=1-101$ and $k=0-200$.

The result of the code is shown here:
\lstinputlisting{NUR_Handin_1_1_output.txt}